\chapter{Discussion}

\section{Introduction}

This chapter brings together the literature, experimental findings and system
evidence to answer the research questions. The discussion separates three
levels of claim. The first concerns discrimination and completion ranking on
the processed Polyvore data. The second concerns the contribution of model
components under controlled ablation. The third concerns the behaviour of the
complete prototype, including candidate search, constraints and fallback.

\section{Relationship to previous compatibility research}

Previous studies have represented an outfit as a sequence, as type-specific
pair relationships, or as a set processed by attention
\parencite{han2017bilstm,vasileva2018typeaware,sarkar2023outfittransformer}.
The present results support the premise shared by these approaches:
compatibility is not adequately represented by isolated visual similarity.
The non-learned cosine and mean-pooled baselines were weaker than the relational
head, while the controlled three-seed comparison favoured the proposed head
over the recurrent, type-aware and set-Transformer adaptations. This evidence
supports the selected project-specific design under the shared protocol. The
comparison models test architecture families under common tensors, frozen
embeddings, labels and training conditions; they are not direct reproductions
of published headline results produced with different experimental protocols.
Under this project's protocol, the small relational head was more effective and
stable than the three controlled alternatives.

The result also depends on frozen SigLIP representations learned from
large-scale image--text pretraining \parencite{zhai2023siglip}. The project
contribution therefore lies in the downstream relational modelling, controlled
evaluation and system integration, while risks inherited from the pretrained
representation remain relevant.

\section{Answer to RQ1: comparative performance}

RQ1 asked how the proposed model compared with non-learned and controlled
learned baselines under consistent protocols. The evidence is positive within
the defined experimental boundary. The selected checkpoint achieved test AUC of 0.8583, with
a 95\% confidence interval of 0.8524--0.8641. It was clearly above
deterministic random scores and mean pairwise cosine similarity. Across three
seeds, its mean AUC exceeded each controlled learned architecture; paired
bootstrap intervals for the AUC differences did not include zero. Its standard
deviation of 0.0009 was also lower than those of the comparison models.

These results indicate that the model learned a repeatable distinction between
observed and corrupted outfits within the processed benchmark
setting. AUC measures the probability that a randomly selected positive
example is ranked above a randomly selected negative example
\parencite{fawcett2006roc}. Same-category replacement reduces trivial
category-based distinctions; compatibility remains defined by the available
data and the published negative-generation protocol.

FITB accuracy of 0.6025 provides a second form of evidence because the model
must select one answer from four alternatives. The result is above the chance
level of 0.25, and its confidence interval does not approach chance. Coverage
limits its interpretation because only 29.50\% of the official questions were
retained. The result applies to
the four supported clothing positions with available embeddings, not to the
complete Polyvore disjoint benchmark. Together, AUC and FITB demonstrate useful
offline discrimination and completion ranking, while the coverage figure
limits the scope of the latter result.

\section{Answer to RQ2: architecture and training choices}

RQ2 asked which design and training choices were supported by the combined
evidence. The strongest structural result concerns explicit pairwise
comparison. Removing the pairwise module reduced mean FITB accuracy by 10.28
percentage points, with a paired 95\% interval from 8.96 to 11.57 points. The
corresponding classification AUC also fell. This supports the view that
compatibility depends on interactions among garments rather than only an
outfit-level average.

The evidence for clothing-position embeddings was weaker. Removing them
reduced mean FITB accuracy by 0.47 percentage points, and the paired interval
ranged from $-0.13$ to 1.06 points. Since the interval crosses zero, the
experiment does not establish a reliable improvement. Its low cost and clear
semantic role provide a practical reason for retention. The evidence is not
strong enough to describe it as necessary. A larger or more position-diverse
evaluation would be needed to resolve the effect.

The selected model was trained with binary cross-entropy on positive and
negative outfits, then used as a scorer for candidate ranking. A scalar
compatibility estimate can order candidates, so this use is coherent. It does
not directly optimise FITB. Under the tested setup, the pairwise ranking-loss
variant did not improve the retained evaluation and was not selected.

Held-out test performance closely matched validation performance, with an AUC
difference of 0.00145. This supports the consistency of checkpoint selection
within the processed split. Training AUC exceeded test
AUC by 0.1116, and validation performance declined after epoch 6, showing that
the model overfit the training distribution. Early stopping controlled the
observed deterioration; it did not eliminate overfitting
\parencite{prechelt1998early}.

Temperature scaling reduced validation negative log-likelihood and expected
calibration error without changing candidate ordering, consistent with its
role as a post-hoc calibration method \parencite{guo2017calibration}. These
improvements apply to the validation distribution. Similarly,
outfit-length weighting was rejected because a small gain in overall AUC
coincided with weaker performance for the lowest-performing length group and a
larger group gap. The decision follows the predefined acceptance gate rather
than a change selected from overall AUC alone.

\section{Answer to RQ3: candidate search and system control}

RQ3 concerned the trade-off between quota-limited retrieval and exact
constrained search over abstract garment states. In the 18-query comparison,
the quota path examined only 3.41\% of combinations on average. This reduced
median CPU time from 0.474 seconds for exact search to 0.028 seconds. At this
lower cost, it recovered the exact highest-scoring result in only 27.78\% of queries and
achieved 24.44\% mean recall of the exact top ten. Mean score loss was positive,
with a bootstrap interval that did not include zero.

The cases were generated deterministically to cover every supported style
twice and every fixed clothing position at least four times. They were not a
random sample of possible users or contexts. The comparison therefore provides
evidence about the implemented search procedures on a bounded coverage set,
not a population-level estimate of retrieval failure.

For the measured candidate space, exact enumeration improved
completeness according to the deployed model objective, while retaining a
median runtime below half a second and a maximum of 3.206 seconds. This supports
its use in the current prototype. It does not show that exact search will
remain practical for an unrestricted catalogue. The number of complete
combinations grows multiplicatively with the candidates in each position, so
later expansion may require pruning or approximate retrieval. Techniques for
large-scale similarity search are well established \parencite{johnson2021similarity}.
Adopting one would introduce a new recall--latency trade-off that should be
measured rather than assumed.

The abstract candidate representation also changed the role of recommendation.
Polyvore did not become a catalogue from which users were instructed to select
archived products. Instead, evidence from the training split was aggregated
into product-independent combinations of clothing position, broad type,
colour, style and clothing range. This matches the interface, which displays a
coloured icon unless an item belongs to the user's local wardrobe. Compared
with the earlier retailer-derived space, the artefact increased visible
type--colour combinations from 122 to 271. This shows broader representational
coverage within the abstract representation.

Search remained subordinate to explicit feasibility rules. The user's
selected item was fixed, weather logic could mask the outer-layer position,
and clothing range was selected manually. All 36 weather--style coverage cases
retained the four logical positions. These checks support functional
completeness within the defined representation.

\section{Achievement of the research objectives}

O1 was met through a critical review that informed a reproducible
secondary-data method and a traceable evidence chain. O2 was met by developing
the frozen-SigLIP compatibility pipeline and comparing its
379,265-parameter head with non-learned, linear, recurrent, type-aware and
attention-based controls, alongside FITB, calibration, ablation, repeated-seed
and overfitting analyses. O3 was met within the prototype boundary by
integrating abstract candidates, exact search, weather and clothing-range
constraints, local wardrobe management and a tested fallback. The interface
uses a short home-page route, soft colour themes and an explicit
menswear--womenswear selector; these are implemented design features rather
than measured usability outcomes. O4 is addressed through the validity and
limitation analysis below, which bounds the findings by FITB coverage, sparse
menswear evidence, source-domain difference and the absence of participant
evaluation.

\section{Validity and confidence in the findings}

Construct validity is limited by the operational definition of compatibility.
The model learned from observed Polyvore outfits and same-category
replacements. These labels are reproducible and cannot capture every reason
why a person might combine garments. FITB is closer to candidate
completion than binary classification, yet it remains an offline choice among
four supplied answers. The resulting evidence therefore concerns ranking
under this dataset-derived definition of compatibility.

Internal validity was strengthened by common splits, frozen features, fixed
comparison settings, validation-based checkpointing and paired analysis. The
local overlap audit found zero shared outfit identifiers, while non-zero item
identifier overlap means that strict product-level disjointness was not
verified and some memorisation risk remains.
The comparison architectures were adaptations, and implementation choices may
favour some models. Three seeds reveal only part of the training variation and
do not characterise all hyperparameter uncertainty.
Bootstrap intervals describe uncertainty conditional on the retained samples
and scoring procedure; they do not correct sampling bias in the original data
\parencite{efron1994bootstrap,bouthillier2021variance}.

External validity is more restricted. The feature-domain classifier separated
Polyvore and Zara embeddings with AUC 1.000. This shows that the two sampled
sources were readily distinguishable in the chosen representation. It does
not quantify how much compatibility performance would change on new retail or
user photographs. The four-position representation excludes dresses and other
multi-position garments, and only 29.50\% of official FITB questions were
evaluable. Menswear evidence is particularly sparse. The scope of the findings
is limited to similar processed data and supported garment structures.

Ecological validity is also untested because no participant study was
conducted. The application demonstrates that its components can be integrated
and that tested actions function; it does not demonstrate usefulness in daily
life. This distinction follows the broader finding that recommender evaluation
requires metrics matched to the intended task and that offline accuracy alone
is incomplete \parencite{herlocker2004evaluation,deldjoo2023review}.

\section{Contribution, benefit and responsible use}

The principal contribution is not a new foundation model. It is an auditable
combination of reused visual representations, a small relational compatibility
head, controlled evaluation and constrained abstract search. The
implementation shows how a resource-limited project can investigate
compatibility without training a large image encoder end to end. Separating
learned scoring from feasibility rules also makes it possible to identify
whether a failure arose from model ranking, candidate availability or a
contextual constraint.

Product-independent output provides a practical design benefit. A
recommendation can describe a broad type and colour without suggesting that
one brand contains the only suitable product. The selected garment remains
visible and fixed, while abstract icons distinguish generated suggestions from
garments saved by the user. Local wardrobe storage and the absence of automatic
gender inference are consistent with the prototype's privacy boundary. These
are implemented properties, not measured user benefits.

A later version could link abstract recommendations to optional product matches
from participating retailers or independent designers. This prospective route
would require verified permissions, current product data and separate
commercial evaluation.

\section{Implications for applied AI practice}

The results suggest that model scale should follow the problem rather than be
treated as the main indication of technical quality. The proposed head trained
faster than the controlled recurrent model and used a similar or smaller number
of trainable parameters than the neural alternatives. Its stronger result came
from an architecture that matched the relational task, especially the explicit
pairwise summary. Once a transferable representation is available, a
focused downstream component can be sufficient for a bounded application.

The development process also shows the value of treating data transformation
as part of the model. The four-position mapping, construction of negatives,
FITB eligibility rules and clothing-range labels determine which relationships
the model can learn and which cases can be evaluated. These choices are not
neutral preparation steps. For example, the exclusion of multi-position
garments directly caused part of the FITB coverage loss, while sparse
menswear-labelled data restricted subgroup conclusions. An applied system
should version its data mapping and report coverage together with
performance, rather than presenting only the final checkpoint.

The same principle applies to candidate generation. The learned score did not
produce useful recommendations until possible garment states were represented,
filtered and searched. The exact-search experiment showed that an apparently
efficient quota could change the result before the model had an opportunity to
score most combinations. In a larger system, retrieval quality would
need its own acceptance criteria. Improving the ranker while leaving an
unmeasured retrieval bottleneck could produce no visible improvement for the
user.

Finally, saved configurations, fixed seeds, hashes, separate
evaluation outputs, explicit constraints, failure-aware loading and a rule
fallback make the path from evidence to interface more inspectable. The neural
representation remains only partly interpretable; these controls reduce
ambiguity about what was trained, selected and deployed. For an MSc Applied AI project,
this system-level traceability is part of technical competence because it
connects an experimental result to reproducible software behaviour.

\section{Prioritised future work}

Future work should address evidence gaps before adding complexity. The first
priority is licensed data expansion. A new dataset should include broader
menswear coverage, documented provenance and explicit structures for dresses
and other one-piece garments. This would require data mapping, new position
occupancy rules, regenerated frozen embeddings, retraining and repetition of
the AUC, FITB, ablation and subgroup evaluation. Dresses should not simply be
relabeled as inner tops because that would incorrectly permit a separate
bottom.

The second priority is target-domain validation. A modest, legally reusable
set of complete outfits from a contemporary target domain should be labelled
under a documented protocol. It could test whether compatibility ordering
transfers beyond Polyvore and whether calibration remains valid. If
performance falls, carefully bounded encoder adaptation or domain-robust
training could be compared with the frozen baseline. Reliable data and
permission review are the main requirements, alongside additional computing
time.

The third priority is participant evaluation after ethics approval. A study
could compare recommendations from the learned model, the rule fallback and a
controlled baseline while measuring perceived usefulness, choice diversity
and trust. Recruitment, informed consent, withdrawal, secure data handling and
a defined analysis plan would be required.

The fourth priority is candidate-search scaling. Exact enumeration should
remain the reference while the current candidate space is bounded. If a larger
catalogue makes latency unacceptable, structured pruning or approximate search
should be evaluated against exact top-$k$ recall, score loss and runtime. This
would preserve a measurable accuracy--efficiency trade-off.
